\section*{Sets and Indices}

\begin{itemize}
    \item $M$ : set of months in the planning horizon, indexed by $m$. For this instance,
    \[
    M=\{1,2,3\}.
    \]
\end{itemize}

\section*{Parameters}

\begin{itemize}
    \item $p_m^{\text{buy}}$ : purchasing price in month $m$ (code: \texttt{purchasing\_prices[m]}).
    \item $p_m^{\text{sell}}$ : selling price in month $m$ (code: \texttt{selling\_prices[m]}).
    \item $C$ : warehouse capacity (code: \texttt{warehouse\_capacity}); here $C=500$.
    \item $I_0$ : initial stock before month $1$ (code: \texttt{initial\_stock}); here $I_0=200$.
\end{itemize}

For this instance,
\[
(p_1^{\text{buy}},p_2^{\text{buy}},p_3^{\text{buy}})=(8,6,9),
\qquad
(p_1^{\text{sell}},p_2^{\text{sell}},p_3^{\text{sell}})=(9,8,10).
\]

\section*{Decision Variables}

For each month $m \in M$:
\begin{itemize}
    \item $x_m \ge 0$ : units purchased at the beginning of month $m$ (code: \texttt{x[m]}).
    \item $s_m \ge 0$ : units sold during month $m$ (code: \texttt{s[m]}).
    \item $I_m \ge 0$ : inventory at the end of month $m$ (code: \texttt{inv[m]}).
\end{itemize}

All decision variables are continuous, exactly as implemented.

\section*{Objective}

Maximize total profit over the quarter:
\[
\max \sum_{m \in M} \left(p_m^{\text{sell}} s_m - p_m^{\text{buy}} x_m\right).
\]

For this instance, this is
\[
\max \; 9s_1+8s_2+10s_3-8x_1-6x_2-9x_3.
\]

\section*{Constraints}

For month $1$:
\[
I_1 = I_0 + x_1 - s_1,
\]
\[
I_0 + x_1 \le C,
\]
\[
I_1 \le C,
\]
\[
s_1 \le I_0 + x_1.
\]

For each month $m \in M \setminus \{1\}$:
\[
I_m = I_{m-1} + x_m - s_m,
\]
\[
I_{m-1} + x_m \le C,
\]
\[
I_m \le C,
\]
\[
s_m \le I_{m-1} + x_m.
\]

Nonnegativity:
\[
x_m \ge 0,\qquad s_m \ge 0,\qquad I_m \ge 0 \qquad \forall m \in M.
\]

\section*{Notes on Implementation}

The implemented model is a linear program solved by PuLP using \texttt{GUROBI\_CMD}. The formulation includes both:
\begin{itemize}
    \item a capacity limit immediately after purchasing, $I_{m-1}+x_m \le C$ (with $I_0$ used for month 1), and
    \item a capacity limit on end-of-month inventory, $I_m \le C$.
\end{itemize}

The sales feasibility constraint is implemented explicitly as
\[
s_m \le I_{m-1}+x_m
\]
(with $I_0$ for month 1), even though it is also implied by inventory balance together with $I_m \ge 0$.

No demand limits, purchase limits, integrality restrictions, or terminal inventory requirements are included in the code, so they are not part of the mathematical model.
